\documentclass{resume} % Use the custom resume.cls style
\usepackage{enumitem} % for item spacing
\usepackage{scalerel} % for inline images

\usepackage[left=0.4 in,top=0.4in,right=0.4 in,bottom=0.4in]{geometry} % Document margins
\newcommand{\tab}[1]{\hspace{.2667\textwidth}\rlap{#1}} 
\newcommand{\itab}[1]{\hspace{0em}\rlap{#1}}
\name{JOHN GIORSHEV} % Your name

\address{
    \scalerel*{\includegraphics{images/location.png}}{|} \href{https://www.google.com/maps/place/Greater+Toronto+Area}{Toronto, ON}\space\space\space\space\space\space
    \scalerel*{\includegraphics{images/email.png}}{|} \href{mailto:johngiorshev@gmail.com}{johngiorshev@gmail.com}\space\space\space\space\space
    \scalerel*{\includegraphics{images/phone.png}}{|} \href{tel:+17058164454}{705-816-4454}\space\space\space\space\space\space
    \scalerel*{\includegraphics{images/linkedin.png}}{|} \href{https://www.linkedin.com/in/johngiorshev/}{johngiorshev}
}  % Your phone number, email, linkedin, and (optional) website

\begin{document}

\begin{rSection}{EMPLOYMENT EXPERIENCE}
    \begin{itemize}
        \item {\bf Software Developer Intern} - SpaceRyde \hfill {Sep 2021 - Sep 2022}
        \begin{itemize}[topsep=-10pt]
            \setlength\itemsep{-0.35em}
            \item[\textbullet] {\bf Telemetry} - Created a module for prioritization, packetization, and streaming of sensors and metrics.\\
            Downstream dashboards using TimescaleDB and Grafana. Data acquisition.
            \item[\textbullet] {\bf Simulator} - Created a 6DOF rigid body simulator for software and controls verification.\\
            Integrated with ROS2, GeographicLib resources, threaded for parallelism, and massively configurable.
            Used to dictate design constraints and purchasing decisions.
            \item[\textbullet] {\bf Hardware Abstractions} - Power distribution unit. STM32 embedded development. MQTT.\\
            Integrated FSMs and QT GUIs over various components for human interface and control.
        \end{itemize}
        \item {\bf Cyber Threat Intelligence Engineering Co-op} - Bell Canada Enterprises \hfill {Summer 2020}
        \begin{itemize}[topsep=-10pt]
            \setlength\itemsep{-0.35em}
            \item[\textbullet] Integrated a MISP server with an Elasticsearch database using RESTful APIs. Elastic Stack. Kafka.\\
            Improved Bell's security information infrastructure. Collaborated using Confluence and Jira.
            \item[\textbullet] Mattermost (slack) notification bot allowing faster team response to latest disseminated exploits.
        \end{itemize}
    \end{itemize}
\end{rSection}

\begin{rSection}{Projects}
    \begin{itemize}
        \item {\bf Key Value} - \href{http://www.keyvalue.ca/}{keyvalue.ca}\\
        Self-hosted full-stack web app for users to store key value pairs. Rust HTTP backend integrated with SQLite. TLS termination using NGINX and Let's Encrypt for secure ingress. Statically served JS, CSS, HTML.\\
        Docker Compose. Applied security best practices (salted password hashes, session tokens, rate limiting).
        \item {\bf Sparse Distributed Representation} - \href{https://github.com/jagprog5/SDR/}{jagprog5/sdr}\\
        C++ sparse math library. Template metaprogramming and specializations for data structures. Integrated with CMake, CI/CD, static analyzers, fuzzing for complete code coverage and profile guided optimization.
        \item {\bf Data Analytics} - \href{https://github.com/jagprog5/resume/blob/main/data-analytics-projects.md}{jagprog5/data-analytics}\\
        Implemented multi-layer perceptrons using PyTorch and C99. Implemented hopfield networks.\\
        Used various sklearn models on Kaggle datasets. Created a shell wrapper of Matplotlib for data visualization. 
        \item {\bf Choose} - \href{https://github.com/jagprog5/choose/}{jagprog5/choose}\\
        Command line interface for lexical analysis and selection dialogs. Regex and unicode support.\\
        Useful for running a line from bash history.
    \end{itemize}
\end{rSection}

\begin{rSection}{Design Teams}
    \begin{itemize}
        \item {\bf QMind} - Queen's AI/ML Software Team \hfill {Sep 2019 - May 2020}\\
        Presented at \href{https://medium.com/qmind-ai/cucai2019-c9d5f848f5c}{CUCAI 2019}. Used OpenCV and Flask to create a computer vision API for retail aisle usage.
        \item {\bf Network Security Team} \hfill {Sep 2018 - May 2019}\\
        Penetration testing on virtual machines with Metasploit. Exploited apps in a sandboxed environment.
\end{itemize}
\end{rSection}

\begin{rSection}{EDUCATION}
    \begin{itemize}
        \item {\bf Computer Engineering} - Queen's University (First Class Honours) \hfill {Graduating May 2023}
    \end{itemize}
\end{rSection}

\end{document}
